\section{Questions suggested by the calculations}
\label{sec:questions}

The ideal volume proposal already has an established measure-theoretic
realization in~\eqref{eq:polyhedralpartition}, and the tree methods described
in Section~\ref{sec:hyperideal} also compute bordered genus-zero volumes.
The following questions isolate further work suggested by the discussion;
they are proposed tasks, not claims that no related results exist.

\begin{question}[Direct hyperideal integration]
Can the fixed-$L$ domains and the density~\eqref{eq:hyperidealdensity} be
integrated directly, with all cell and symmetry factors, to recover
$2\pi^2+\frac12\sum_iL_i^2$? Which change of variables makes the
cancellation or polynomiality visible?
\end{question}

\begin{question}[A usable measure on general polyhedral cells]
Can~\eqref{eq:pfaffiandensity} be expressed in dihedral-angle variables in
a way that makes higher-dimensional cell integrals computationally useful?
In the ideal case the measure is particularly simple. What replaces that
simplicity for fixed hyperideal boundary lengths or compact cone openings?
\end{question}

\begin{question}[Gluing and topology]
Can a polyhedral description be combined with Fenchel--Nielsen gluing to
give a practical calculation for arbitrary $(g,n)$? A proposed formula
must specify both the cutting multiplicity and the full mapping-class
quotient; integrating over all pants lengths and twist periods alone
overcounts moduli space.
\end{question}

\begin{question}[Curvature corrections in period coordinates]
Can the correction class $\Xi_\mu$ in~\eqref{eq:Xiclass} be represented
by an explicit first curvature correction to the WP form in
the actual finite-area conformal coordinates? Equation~\eqref{eq:Frickefirstcorrection}
answers the local four-point calculation in a blown-up character chart,
whose conversion to conformal coordinates is still required.
How does the local density behave
in a joint limit toward a balanced collision wall, where the
final-chamber expansion is not uniform in the weights?
\end{question}

\begin{question}[Beyond the regular quantum degeneration]
Can Proposition~\ref{prop:regularquantumAiry} be extended uniformly
over compact sets of nondegenerate Euclidean tetrahedral shapes,
with a geometric description of the required scalar normalization?
Can a geometric state
integral in that regime produce the Thurston--Deligne--Mostow measure
on the boundary-shape space? The model~\eqref{eq:modelpartition}
specifies a classical comparison, not an identification with an
existing quantum invariant. The exact Liouville
identity~\eqref{eq:exactLiouvillecurvature} and the torsion
limit~\eqref{eq:torsionDMlimit} identify two precise boundary
constructions; the auxiliary-tetrahedron obstruction and the
different regimes in~\eqref{eq:areaLiouvillecharge} explain
what remains to connect them to that joint asymptotic.
The finite one-screening fusion calculation in
Section~\ref{sec:screenedfusion} computes the finite period and
residue-sector comparison, with a logarithmic limit at balanced
weights. Its direct infinite-spectral derivation uses the
continuation and tail-control input~\eqref{eq:screeningtailinput};
Section~\ref{sec:normalizedsixj} fixes the quantum-group
normalization, proves the rank-two obstruction, and computes
the classical rank-three fold in the regular compact model.
Section~\ref{sec:quantumAiry} corrects the integration-coordinate
dictionary, removes the external parameter poles, and proves the
uniform regular-family Airy formula with its first correction.
General shape dependence, global gluing, and identification of
the resulting measure remain to be established.
\end{question}

\begin{question}[Uniform quantum Dehn filling]
For a knot complement admitting the positive geometric optimization
of Section~\ref{sec:dehnresponse}, can an exact quantum surgery
representation be controlled uniformly in a complex neighborhood
of complete peripheral holonomy, with a nonzero geometric contribution
and bounds on all competing terms? Proposition~\ref{prop:dehncriticalvalue}
already fixes every classical long-slope coefficient, and
\eqref{eq:dehngaussiantorsion} specifies the torsion check. The further
issue is the analytic input of
Proposition~\ref{prop:dehnquantumconditional}, especially in a joint
limit where the slope varies with the quantum level.
\end{question}


\begin{question}[Certified realization and quantitative angle bounds]
For the planar problem of Section~\ref{sec:ipolybridge}, can positive
feasibility, initial centering, inexact solves, and geometric error
certification be combined into a complete complexity theorem with
explicit input and precision dependence? For a three-manifold angle
polytope whose maximum is interior, which additional geometric hypotheses
bound its smallest angle in terms of the combinatorial input and a
quantitative feasible margin? The barrier iteration bound is available;
the planar Laplacian sign structure does not extend automatically to
the full three-dimensional Hessian.
\end{question}

\begin{question}[A continuum comparison of second variations]
Can a family of refined polyhedral functionals, with specified geometric
convergence and normalization, have a second-variation limit related to
a WP potential? The exact cotangent identity, Schur complement, and
Gaussian determinant in Section~\ref{sec:ipolybridge} provide discrete
objects for such a comparison. Control of their limits and an explicit
identification of the target metric are still required.
\end{question}

\begin{question}[Continuation of a coefficient Laplace transform]
Does $D_C(s)$ from Section~\ref{sec:coefficientdistributions} admit a
meromorphic continuation beyond the proved half-plane $\Re s>1$?
A proof must control coefficient extraction in the continued
character-valued expression, since it is not bounded for the
uniform norm on unitary characters.
\end{question}

\begin{question}[Control of the globally filtered remainder]
Proposition~\ref{prop:globalonecrossing} gives a bounded global
cancellation of every one-crossing primitive and all its iterates.
Can the signed residual $\mathcal R_X$ in~\eqref{eq:globalresidual}
be estimated with constants controlled by the geometry of $X$?
The exact obstruction~\eqref{eq:globalobstruction} rules out a
perfect fixed linear selector; length-dependent cancellation or
the nonlinear diagonal extraction in~\eqref{eq:globaldiagonal}
retains information that such a selector discards.
\end{question}

\begin{question}[Effective selection of connected supports]
Can the signed coefficient relations produce the smoothed estimate
\eqref{eq:countsmoothhypothesis} with specified $\eta,m$ and geometric
constants? The elementary lattice remainder, the stable finite support
projection~\eqref{eq:countsupportfourier}, and the constructive isolation
estimate are available. The global one-crossing filter and the
length-valued Euler product now provide concrete aggregate objects.
An estimate for the higher-intersection residual or for the nonlinear
first-atom extraction, strong enough to yield a counting saving,
remains to be proved.
\end{question}

\begin{question}[From weighted curve counts to mixing]
Can the filling-pair coefficient construction of
Section~\ref{sec:pairconverse} produce a norm-controlled estimate
for the weighted orbit kernel~\eqref{eq:twopointkernel} without
assuming mixing? Proposition~\ref{prop:twopointunfolding} gives the
exact unfolding, and Theorem~\ref{thm:twopointfullrate} preserves
the counting exponent when derivative endpoint weights are controlled. The independent
two-endpoint estimate, its Jacobian weights, and the noncompact tail
remain to be established for that principal-stratum kernel; a scalar
area or length count is insufficient. Section~\ref{sec:genus2markedtorus}
settles the separate affine return estimate on the nonsquare $(2,2)$
stratum, and Section~\ref{sec:coefficientobservables} gives selected
primitive-arc coefficient observables with uniform geometric norms.
Can their class be enlarged to control the full two-endpoint kernel,
with quantitative stability of the nonlinear coefficient extraction,
and can the calculation be extended beyond the marked-torus locus?
\end{question}

\begin{question}[The sharp norm-ball remainder]
Can the qualitative $o_X(L\log L)$ remainder of
Proposition~\ref{prop:strictconvexlittleo} be made effective through a
decay modulus satisfying~\eqref{eq:primitiveDini}, or through cancellation
in~\eqref{eq:primitiveexactnormalizedmobius}? The inverse-flatness
formula~\eqref{eq:inverseflatmicroscopic} quantifies a local lower
obstruction, while~\eqref{eq:dyadicerrorinverse} gives a test independent
of the leading coefficient. Determining whether the primitive remainder
is $O_X(L)$, or has an unbounded slowly varying factor, remains open here.
Section~\ref{sec:arithmeticgeometry} makes the two inputs explicit:
$\Delta_B(u)=O(u/\log u)$ together with any fixed Mertens power
saving gives $O_X(L)$, whereas an integrable geometric modulus and
the prime number theorem already give $o_X(L)$.
The arithmetic and flatness hypotheses must be specified for each surface;
finite-range numerical observations do not decide them.
\end{question}

\begin{question}[Beyond closure operations]
Which positive constant series have a geometric partition or a useful
functional equation behind them? The figure-eight kernel
\eqref{eq:universalkernel} and the twist kernel~\eqref{eq:Kkernel} provide
two concrete models. Their different index sets suggest separating
global curve identities from identities for a distinguished twist orbit.
\end{question}

\begin{question}[A boundary phase-space formulation]
Can the endpoint-and-direction matching relations in
Section~\ref{sec:matching} be organized into an integration scheme for
general filling curves? Orthogeodesic lengths give the canonical local
geometry, while the endpoint phase variables retain the missing gluing
information. A useful result would control both the sum over itineraries
and the integration over those variables.
\end{question}

These questions concern concrete integrals or pointwise identities. The
distinctions established above---fixed parameters, geometric chambers,
correct curve types, and specified measures---are part of their statements,
rather than technicalities to be added after a formal calculation.

